\section*{अध्याय \ref{ch:g12:synthesis-strategy} --- संश्लेषण की रणनीति और हरित रसायन}

\begin{solution}{exo:g12:synthesis-strategy:1}
जलयोजन: $46.0/(28.0 + 18.0) = \qty{100}{\percent}$। किण्वन:
$2 \times 46.0/180.0 = \qty{51.1}{\percent}$।
\end{solution}

\begin{solution}{exo:g12:synthesis-strategy:2}
दोनों समूह, ग्लाइसीन के \omterm{def:g11:functional-groups:amine}{ऐमीन} पर \ce{Z} और ऐलानीन के अम्ल पर मेथिल \omterm{def:g11:functional-groups:carboxylic-acid}{एस्टर},
चरण 1 में लगाए जाते हैं और चरण 3 में हटाए जाते हैं।
\end{solution}

\begin{solution}{exo:g12:synthesis-strategy:3}
सिद्धांत 5: अधिक सुरक्षित \omterm{def:g6:solutions-and-solubility:solute}{विलायक} और अभिक्रिया परिस्थितियाँ इस्तेमाल कीजिए।
\end{solution}

\begin{solution}{exo:g12:synthesis-strategy:4}
एथेनॉल की अधिकता $Q < K$ कर देती है, इसलिए अभिक्रिया और आगे
बढ़ती है। \omterm{def:g12:catalysis:catalyst}{उत्प्रेरक} दोनों दिशाओं को समान रूप से तेज़ करता है: वही \omterm{def:g10:reaction-progress-table:system}{अंतिम
अवस्था} प्राप्त होती है, बस जल्दी।
\end{solution}

\begin{solution}{exo:g12:synthesis-strategy:5}
हाँ: केवल \omterm{def:g11:functional-groups:carbonyl}{कीटोन} अपचयित होता है, एक द्वितीयक \omterm{def:g11:functional-groups:alcohol}{ऐल्कोहॉल} समूह में;
\omterm{def:g11:functional-groups:carboxylic-acid}{एस्टर} अपरिवर्तित रहता है।
\end{solution}

\begin{solution}{exo:g12:synthesis-strategy:6}
अधिक \omterm{def:g11:functional-groups:carboxylic-acid}{एस्टर}: किसी एक \omterm{def:g7:chemical-reactions:reactant}{अभिकारक} की अधिकता, या बनते ही पानी हटाना (जल-पाश)।
अधिक तेज़: \omterm{def:g10:synthesis-yield:reflux}{पश्चवाह तापन} कीजिए, अम्ल \omterm{def:g12:catalysis:catalyst}{उत्प्रेरक} मिलाइए।
\end{solution}

\begin{solution}{exo:g12:synthesis-strategy:7}
\ce{NaBH4} के साथ:
\[ \ce{CH3-CH(OH)-CH2-CH2-CO-O-CH3} . \]
\ce{LiAlH4} के साथ:
\[ \ce{CH3-CH(OH)-CH2-CH2-CH2OH} \]
(पेंटेन-1,4-डाइऑल), साथ में \omterm{def:g11:functional-groups:carboxylic-acid}{एस्टर} के दूसरे भाग से मेथेनॉल।
\end{solution}

\begin{solution}{exo:g12:synthesis-strategy:8}
योगज (\qty{100}{\percent}), एस्टरीकरण
(\qty{83}{\percent}) और आइबुप्रोफ़ेन तक तीन-चरण वाला रास्ता
(\qty{77}{\percent})। $77.4/40.0 = 1.9$: लगभग दोगुना अच्छा।
\end{solution}

\begin{solution}{exo:g12:synthesis-strategy:9}
$180.0/(138.0 + 102.0) = \qty{75.0}{\percent}$। एथेनोइक अम्ल:
प्रति किलोग्राम ऐस्पिरिन $60.0/180.0 = \qty{0.333}{kg}$।
\end{solution}

\begin{solution}{exo:g12:synthesis-strategy:10}
तापन अभिक्रिया को तेज़ करता है; संघनित्र वाष्पों को फ़्लास्क में लौटा देता है,
ताकि कोई \omterm{def:g7:chemical-reactions:reactant}{अभिकारक} या \omterm{def:g6:solutions-and-solubility:solute}{विलायक} नष्ट न हो। तापन दर सुधारता है,
अंतिम \omterm{def:g10:synthesis-yield:yield}{लब्धि} नहीं।
\end{solution}

\begin{solution}{exo:g12:synthesis-strategy:11}
छह चरण बनाम तीन। सिद्धांत: अपशिष्ट रोकना (1), \omterm{def:g12:synthesis-strategy:atom-economy}{परमाणु मितव्ययिता} अधिकतम
करना (2), \omterm{def:g12:catalysis:catalyst}{उत्प्रेरक} इस्तेमाल करना (9); साथ ही कम चरण, कम ऊर्जा (6)।
\end{solution}

\begin{solution}{exo:g12:synthesis-strategy:12}
$0.80 \times 0.70 \times 0.90 = 0.504$: \qty{50.4}{\percent}। प्रारंभिक
सामग्री: $0.10/0.504 = \qty{0.20}{mol}$।
\end{solution}

\begin{solution}{exo:g12:synthesis-strategy:13}
ऐलानीन के \omterm{def:g11:functional-groups:amine}{ऐमीन} (\ce{Z}) और ग्लाइसीन के अम्ल (मेथिल \omterm{def:g11:functional-groups:carboxylic-acid}{एस्टर}) की रक्षा कीजिए;
ऐलानीन के मुक्त अम्ल और ग्लाइसीन के मुक्त \omterm{def:g11:functional-groups:amine}{ऐमीन} के बीच \omterm{def:g11:functional-groups:amine}{ऐमाइड} आबंध
बनाइए; दोनों \omterm{def:g12:synthesis-strategy:protecting-group}{रक्षी समूह} हटाइए। तीन अवस्थाएँ, अर्थात पाँच
अभिक्रियाएँ, यदि हर रक्षण और विरक्षण गिना जाए।
\end{solution}

\begin{solution}{exo:g12:synthesis-strategy:14}
\omterm{def:g12:catalysis:catalyst}{उत्प्रेरक} सिद्धांत 9 लागू करता है। पर \qty{250}{\celsius} सिद्धांत 6
के विरुद्ध है, बेंज़ीन सिद्धांत 3 और 5 के विरुद्ध, और
\qty{35}{\percent} की \omterm{def:g12:synthesis-strategy:atom-economy}{परमाणु मितव्ययिता} सिद्धांत 2 के विरुद्ध: प्रति किलोग्राम उत्पाद
$65/35 = \qty{1.9}{kg}$ अपशिष्ट। एक हरित विशेषता किसी प्रक्रम को हरित
नहीं बना देती।
\end{solution}

\begin{solution}{exo:g12:synthesis-strategy:15}
% ledger: crit:CO2-T, crit:CO2-p
\qty{40}{\celsius} और \qty{100}{bar} कार्बन डाइऑक्साइड के क्रांतिक बिंदु,
\qty{31}{\celsius} और \qty{73.8}{bar}, से ऊपर हैं। दाब
वायुमंडलीय दाब का सौ गुना है: पात्र मोटे इस्पात का होता है। कार्बन
डाइऑक्साइड न विषैली है न ज्वलनशील, और वह बिना अवशेष छोड़े
फिर से गैस बन जाती है, और दोबारा इस्तेमाल की जा सकती है; कोई क्लोरीनयुक्त \omterm{def:g6:solutions-and-solubility:solute}{विलायक} बाहर नहीं निकलता।
\end{solution}

\begin{solution}{pb:g12:synthesis-strategy:1}
% ledger: ibu:bhc-ae-recovered
\textbf{1.} C: $10 + 4 + 1 = 15 = 13 + 2$; H: $14 + 6 + 2 = 22 = 18 + 4$;
O: $3 + 1 = 4 = 2 + 2$।

\textbf{2.} एथेनोइक अम्ल, \ce{C2H4O2}।

\textbf{3.} नाइट्रोजन, क्लोरीन और सोडियम: आइबुप्रोफ़ेन में केवल कार्बन,
हाइड्रोजन और ऑक्सीजन हैं।

\textbf{4.} नया रास्ता: सिद्धांत 9, रससमीकरणमितीय \omterm{def:g7:identifying-substances:characteristic-test}{अभिकर्मकों} के बजाय
\omterm{def:g12:catalysis:catalyst}{उत्प्रेरक}।

\textbf{5.} $13 \times 12.0 + 18 \times 1.0 + 2 \times 16.0 =
\qty{206.0}{g/mol}$।

\textbf{6.} $134.0 + 102.0 + 122.5 + 68.0 + 19.0 + 33.0 + 36.0 =
\qty{514.5}{g/mol}$।

\textbf{7.} $206.0/514.5 = \qty{40.0}{\percent}$।

\textbf{8.} $134.0 + 102.0 + 2.0 + 28.0 = \qty{266.0}{g/mol}$।

\textbf{9.} $206.0/266.0 = \qty{77.4}{\percent}$।

\textbf{10.} काग़ज़ पर $(206.0 + 60.0)/266.0 = \qty{100}{\percent}$; व्यवहार में
\qty{99}{\percent} से अधिक।

\textbf{11.} $514.5/206.0 = \qty{2.50}{t}$।

\textbf{12.} $2.50 - 1 = \qty{1.50}{t}$।

\textbf{13.} $266.0/206.0 = \qty{1.29}{t}$।

\textbf{14.} $60.0/206.0 = \qty{0.29}{t}$।

\textbf{15.} $1.50 - 0.29 = \qty{1.21}{t}$।

\textbf{16.} $0.90^6 = 0.53$: \qty{53}{\percent}।

\textbf{17.} $0.90^3 = 0.73$: \qty{73}{\percent}।

\textbf{18.} हर चरण में \omterm{def:g6:solutions-and-solubility:solute}{विलायक}, पृथक्करण और शोधन लगते हैं, और
कुछ उत्पाद नष्ट होता है: कम चरण, ऐसी कम हानियाँ।

\textbf{19.} लगभग कुछ नहीं: उसका एकमात्र उप-उत्पाद इस्तेमाल हो जाता है।

\textbf{20.} तीन-चरण वाला रास्ता प्रति टन आइबुप्रोफ़ेन लगभग \qty{1.5}{t} अपशिष्ट
बचाता है (\qty{1.2}{t}, भले ही उसका एथेनोइक अम्ल फेंक दिया
जाए)।
\end{solution}
